\section{Sparse obstacles and local propagation}\label{sec:geometry}

The return comparison requires an exceptional set smaller than
$\beta^{-5/2}$ and spatial propagation sharper than a ballistic bound.
We obtain both from a six-point sparsity estimate and a displacement
estimate for isolated exposed strands.  The former is an extension of
the one-point argument in \cite[Section~4]{BlochLaw}; the latter uses a
direct recoloring of the exchange picture.  We then fix a good exposure
and derive the heat-kernel and flow estimates used in the comparison.

All site distances and balls in this section use the supremum norm.
A site is \emph{blocked} if none of its slots is a hole.  In the extended
hole graph, blocked sites occur only inside $\Lambda$.

\subsection{Geometric inputs and the snapshot estimate}

We first recall the geometric hypotheses under which the available
slots support uniform diffusion and flows.  Fix the constant
$M_*\ge48$ used in \cite[Section~3]{BlochLaw}; it is chosen there from
\cite[Lemma~5.5]{Ordering}.  A time is \emph{regular at radius $R_0$
about $x$} if:
\begin{enumerate}
\item every $B_{M_*}(z)$ with $z\in B_{11R_0}(x)$ contains fewer than
six blocked sites;
\item for each blocked $z$ with $b=|z-x|_\infty>10R_0$, its blocked
component, with component steps of supremum length at most eight, has
radius about $z$ at most $(1+b)^{1/300}$.
\end{enumerate}
We call the first condition \emph{local regularity}.

\begin{proposition}[Imported geometric estimates]\label{prop:geometryinputs}
For all sufficiently large $\beta$, every allowed pin set $D$, every
deterministic time, every anchor $x$, and every $R_0$ above a fixed
threshold, the following bounds hold.  Put $p=(5/\beta)^{10/11}$.
\begin{enumerate}
\item The probability of a nonregular time under $\PP_D$ is at most
\begin{equation}\label{eq:regularityprob}
 C R_0^3p^6+C e^{-cR_0^\xi},
\end{equation}
for some fixed $\xi>0$.  Failure of the second regularity condition
alone has probability at most $C e^{-cR_0^\xi}$.
\item At every locally regular time, every nonnegative vector on the
hole graph supported above $B_{R_0}(x)$ satisfies
\begin{equation}\label{eq:nash}
 \norm{v}_2^{2+4/3}\le C\D_t(v)\norm{v}_1^{4/3}.
\end{equation}
There is an available slot within bounded distance of every site in
$B_{2R_0}(x)$.  Available slots over two sites in this ball at distance
$b$ can be joined by a path of length at most $C(1+b)$, staying within
distance $C(1+b)$ of the first site and inside $B_{3R_0}(x)$.
\item At every regular time, each available slot $a$ above $B_{R_0}(x)$
admits a unit flow $J_a$ to infinity on the extended hole graph.
For $b$ the distance between its source site and an edge,
\begin{equation}\label{eq:flowenvelopes}
 |J_a(e)|\le C(1+b)^{-199/100},
 \qquad |J_a(e)|\le C(1+b)^{-2}\quad(b\le R_0).
\end{equation}
\end{enumerate}
All constants are uniform in the volume, pins, anchor, and deterministic
time.
\end{proposition}

The probability estimates are
\cite[Proposition~5.4 and Theorem~7.5]{Ordering}, with their separate
local and distant-component bounds, as applied in
\cite[Proposition~4.3]{BlochLaw}.  To check the time-grid hypothesis of the
pin theorem, include zero, the queried time, and every extra pin time
in a finite grid and refine until every gap is at most $p^{12}$.
Then $\beta=5p^{-11/10}$, and the theorem is uniform in the grid.

The Nash inequality is \cite[Proposition~5.6]{Ordering}.  Its proof
uses only local regularity: alternatively fill all sufficiently remote
sites with holes to obtain full regularity without changing the energy
of a vector supported in $B_{R_0}(x)$.  The path assertion follows from
\cite[Lemma~5.5]{Ordering}.  Every site in $B_{2R_0+1}(x)$ has
an unblocked nearest neighbor, since six blocked neighbors would violate
local regularity.
Along a nearest-neighbor site route replace blocked vertices by such
neighbors and use the lemma's bounded paths between successive surviving
vertices.  All detours have bounded size.  Slots over the same site are
connected through an unblocked neighbor.  For $R_0$ above a fixed
threshold the resulting paths stay in $B_{3R_0}(x)$.
The global flow envelope is \cite[Proposition~6.2]{Ordering}; the
inverse-square improvement is proved in
\cite[Lemma~3.6]{BlochLaw}.  That proof uses only regularity at
the stated radius and a fixed lower threshold, so the near envelope is
available for arbitrary $R_0$ here.

We next strengthen the pin theorem to its natural diffusive density,
for the finitely many points needed below.  Write
$\lambda=\beta^{-3/2}$.

\begin{proposition}[Uniform snapshot sparsity]\label{prop:snapshot}
For every allowed pin set $D$ and every set $T$ of distinct slot points
at one deterministic time, with $1\le |T|\le6$,
\begin{equation}\label{eq:snapshot}
 \PP_D(T\subset\mathcal B_D)\le C\lambda^{|T|}.
\end{equation}
The bound is uniform in the finite cube, the pin set, the time, and the
points, including distinct slots at the same site.
\end{proposition}

\begin{proof}
If $T$ meets $D$, the probability is zero.  Otherwise set $m=|T|$ and
$P'=D\cup T$.  By \eqref{eq:simultaneouspin}, it suffices to prove
$r_T\le C\lambda$.  We adapt the one-turn energy argument of
\cite[Section~4]{BlochLaw}, retaining all $m$ sources.

Fix the exposure under $\PP_{P'}$ and use the time of $T$ as the cut.
Let $U$ be the fresh one-period kernel killing only on the inner
boundary, and put $Q=(U+U^*)/2$ and $\eta=\ind_T$.  Let $f=1$ on $T$;
off $T$, let $f$ be the probability of hitting $T$ at a future cut
before boundary killing.  Finite-horizon approximation gives
$0\le f\le1$ and $U^*f=f$ off $T$, even if eventual killing is not
almost sure.  For $g=f-\eta$, put $e_g=\ip{g}{(I-Q)g}\ge0$.
Removing all pins of $D$ except the boundary increases the return
probability.  Harmonicity off $T$ and expansion of the quadratic form
therefore give
\begin{align}
 m r_T
 &\le\EE_{P'}\bigl[m-\ip{f}{(I-Q)f}\bigr]\notag\\
 &=\EE_{P'}\bigl[
      \ip{\eta}{Q\eta}+\ip{(U+U^*)\eta}{g}-e_g\bigr].
 \label{eq:snapshotenergy}
\end{align}
This is also the multipoint hitting-energy identity in
\cite[Lemma~4.5]{Ordering}.

Here are the propagation estimates used to bound the two positive
terms.  For each $j\in T$, let $x_j$ be its site and choose radius
$R_0=\lceil a\beta\rceil$, with $a$ a sufficiently large fixed
constant.  Let $H_j$ be the event that the total nonregular time at this
radius about $x_j$ is at most $\beta/100$.
Equation~\eqref{eq:regularityprob} and Markov's inequality give
\begin{equation}\label{eq:bootstrapbad}
 \PP_{P'}(H_j^c)+
 \sup_t\PP_{P'}(t\text{ is nonregular about }x_j)
       \le C\beta^{-27/11}=o(\lambda).
\end{equation}
The localization, diffusion, and spread-source estimates in
\cite[Lemma~4.4 and Propositions~4.5 and~3.7]{BlochLaw}, with one turn, have the following
consequences.  Replacing a propagation from $j$ by the propagation killed
on exit from $B_{R_0}(x_j)$ loses expected mass at most $Ce^{-c\beta}$.
On $H_j$, a localized propagation of duration at least $\beta/3$ has
$\ell^2$ norm at most $C\beta^{-3/4}$; localized full-period entries are at most
$C\lambda$.  At a regular endpoint, a nonnegative localized mass $\mu$
of mass at most one and norm at most $C\beta^{-3/4}$ satisfies
\begin{equation}\label{eq:bootstrapspread}
 |\ip{\mu}{Y}|\le
                 \bigl(C\beta^{-1/2}\D_t(Y)\bigr)^{1/2}
\end{equation}
for every finitely supported $Y$.  These estimates apply with pin law
$P'$: for the source $j$, all other points of $P'$ are allowed extra
pins.  They impose no separation on $T$.

Since there are at most $m^2$ return entries, localization, diffusion,
and \eqref{eq:bootstrapbad} imply
$\EE_{P'}\ip{\eta}{Q\eta}\le C\lambda$.
For the cross term, consider one pairing $\ip{U_d\ind_{\{j\}}}{g}$,
where $U_d=U$ or $U^*$.  Split it at elapsed time
$u\in I=[\beta/3,2\beta/3]$.  Localize only the propagated starting
mass, and let $Y(u)$ be the adjoint suffix applied to $g$.
Both the mass and $\norm{Y(u)}_\infty$ are at most one.
Discarding the localization loss, $H_j^c$, and nonregular split slices
therefore changes the expected pairing by $o(\lambda)$.
On retained slices use \eqref{eq:bootstrapspread}; then average over
$I$.  Lemma~\ref{lem:kernelenergy} and Cauchy--Schwarz give, at each
fixed exposure, the retained average bound
\[
 \frac{C\beta^{-1/4}}{|I|}
         \int_I\sqrt{\D_{t(u)}(Y(u))}\,du
       \le C\beta^{-3/4}\sqrt{e_g}.
\]
In this application the energy uses boundary-only killing, as does
$U$, and its quadratic form equals that of $Q$.
There are $2m\le12$ such pairings.  Substituting in
\eqref{eq:snapshotenergy} and using
\[
 C\beta^{-3/4}\sqrt{e_g}-e_g\le C'\beta^{-3/2}
\]
proves $r_T\le C\lambda$.  Equation~\eqref{eq:simultaneouspin}
now proves \eqref{eq:snapshot}.
\end{proof}

\subsection{Short blocks and isolated strands}

The density $\lambda$ in Proposition~\ref{prop:snapshot}, in place of
the coarser $p$ in \eqref{eq:regularityprob}, makes six-obstacle events
sufficiently rare on the larger spatial scales introduced next.
The snapshot estimate will give at most five obstacles in a mesoscopic
neighborhood.  To compare an isolated moving obstacle with a smooth
profile, we also need its displacement to be smaller than the block
length.  We prove that displacement bound before assembling the good
event, so that no conditioning on that event enters its proof.

From now on $i$ lies at time zero just after the seam, over site $x$,
and $P=E\cup\{i\}$.  Choose
\begin{equation}\label{eq:scales}
 D_*=\sqrt\beta,\qquad L=\beta^\delta,\qquad\delta=\frac1{200},
 \qquad R=\beta^{3/2+\delta/4},\qquad
 R_*=\lceil\beta^{5/3}\rceil.
\end{equation}
We assume $B_{50R_*}(x)$ is interior to $\Lambda$ and contains no site
above which $E$ has a pin.  All subsequent bounds are uniform under this
assumption.  We call an error \emph{rapidly small} if, for every fixed
$j$, its absolute value is at most $C_j\beta^{-j}$ with this uniformity.

Partition each half-period into equal blocks of length between $L$ and
$2L$.  Subdivide the blocks into bins of length between one and three.
Use the same partitions in reverse.  The seam is assigned to the last
forward block and bin and to the first reverse block and bin.  An
exposed \emph{strand} in a block is the trajectory of one exposed slot
at its starting slice in the chosen direction.  This definition counts
different strands of the same cycle separately.

\begin{lemma}[Domination of an isolated strand]\label{lem:isolatedpath}
Fix an allowed pin set $P$, a deterministic slot $v$, and an interval
$[s,t]$ of positive length less than $\beta$ containing no seam in its
interior.  Follow the picture's strand $X$ from $v$ at time $s$ to
time $t$.  Let $\mathcal I$ be the event that this strand is exposed
and, throughout the interval, no other exposed slot lies at its site
or a neighboring site, on both sides of every jump.  For every
measurable path event $\mathcal A$ consisting of paths that stay in
the interior,
\begin{equation}\label{eq:isolateddomination}
 \PP_P\bigl(\mathcal I\cap\{X\in\mathcal A\}\bigr)
       \le2\EE_{\mathrm b}\ind_{\{X\in\mathcal A\}}.
\end{equation}
On the right, $X$ is the base tagged walk, with rate $1/2$ along each
slot edge; on interior paths its law is the free walk law.  Any seam
at an endpoint is excluded: use the post-seam side at $s$ and the
pre-seam side at $t$.  The same bound holds when both paths are traced
backward from a prescribed slot at $t$ to $s$, with these same
chronological sides.
\end{lemma}

\begin{proof}
We first impose that the distinguished exposed cycle is colored down.
Conditionally on a picture in $\mathcal I$, this costs exactly a factor
$1/2$.  We prove that the resulting joint path law is bounded by the
base tagged-walk law.

Fix a path $X$ starting at $v$, with successive slots joined by slot
edges and jump times $s<\tau_1<\cdots<\tau_j<t$.
Let $A_X$ consist of the edge--time pairs $(e,u)$ with $s<u<t$,
$u$ not a jump time, and $e$ incident to $X(u)$; let $\mu$ be the base
Poisson intensity measure.  The path requires its $j$ jump marks and
the absence of all other marks in $A_X$.  Poisson disintegration gives
the density
\[
 q(X)=2^{-j}\exp\{-\mu(A_X)\}
\]
with respect to counting measure on such slot lists and Lebesgue
measure on ordered jump times.  The complementary clocks and seams
retain their independent base laws.  Denote those variables by $\zeta$,
and let $N_X(\zeta)$ count the admissible color histories for which
$\mathcal I$ holds and the distinguished cycle is down.

Recolor only the segment $X$ up.  Every other slot at its site or a
neighboring site is up, since isolation makes those slots unexposed.
At a prescribed jump the colors change from
$(\downarrow,\uparrow)$ to $(\uparrow,\downarrow)$; after recoloring,
both slots are up on both sides, so the mark can be deleted.  Fill
$A_X$ with arbitrary additional marks.  Each new mark joins two up
slots, so the recolored history remains valid.  Its only color
discontinuities are the prescribed $\downarrow\!\to\!\uparrow$ flip
at $(v,s)$ and $\uparrow\!\to\!\downarrow$ flip at $(X(t),t)$.
All pins remain up, because the original segment belongs to an
exposed cycle.  Figure~\ref{fig:surgery} depicts the deleted jumps,
the arbitrary up--up fillings, and the oriented endpoint flips.

For fixed $X$, $\zeta$, and filling, the map on color histories is
injective: recoloring the specified segment down recovers the original
history.  In the filled picture, a cycle containing a prescribed
oriented flip has at most one compatible coloring.  Every other cycle
has at most its usual one choice if pinned and two choices if unpinned.
Writing $w_P=2^{n_P}$ and $\omega'$ for the filled picture, we therefore
have $N_X(\zeta)\le w_P(\omega')$ for every filling.

For fixed $X$, the set $A_X$ is deterministic.  Independent Poisson
filling on $A_X$, together with the complementary clocks and seams,
has exactly the unconditioned base law.  Averaging yields
\[
 \EE_\zeta N_X(\zeta)\le\EE_{\mathrm b}w_P=Z_P.
\]
The down-strand path density is
$q(X)\EE_\zeta N_X(\zeta)/Z_P$, hence at most $q(X)$.
Summing over jump lists and integrating their ordered times proves
the assertion for every measurable $\mathcal A$; finite volume gives
finitely many marks almost surely.  Removing the down-color condition
supplies the factor two.  At seam endpoints the flips are inserted on
the specified sides of the unchanged seam mapping.  Reversing time
gives the reverse assertion, with the flip directions reversed.
\end{proof}

\begin{figure}[tbp]
\centering
\begin{tikzpicture}[x=0.86cm,y=0.88cm,>=Stealth,
 every node/.style={font=\small}]
 \definecolor{downpath}{RGB}{165,55,50}
 \definecolor{upfill}{RGB}{25,108,105}
 \begin{scope}
  \node[font=\small\bfseries] at (1.5,4.5) {Fixed isolated down strand};
  \foreach \x in {0,1,2,3}
   \draw[gray!40] (\x,0.25)--(\x,3.8);
  \draw[->,gray] (-0.65,0.35)--(-0.65,3.7) node[above] {$t$};
  \draw[downpath,very thick] (1,0.35)--(1,1.15)--(2,1.15)--(2,2.55)--(1,2.55)--(1,3.7);
  \draw[downpath,->,thick] (1.2,1.15)--(1.85,1.15);
  \draw[downpath,->,thick] (1.8,2.55)--(1.15,2.55);
  \fill[downpath] (1,0.35) circle (2pt);
  \fill[downpath] (1,3.7) circle (2pt);
  \node[left] at (0,0.35) {$t_0$};
  \node[left] at (0,3.7) {$t_1$};
  \node[right,align=left,text=downpath] at (2.2,1.6) {$X$};
  \node[align=center] at (1.5,-0.35) {The incident holding intervals\\contain no exchange marks.};
 \end{scope}
 \begin{scope}[xshift=6.8cm]
  \node[font=\small\bfseries] at (1.5,4.5) {Recolor, delete, and fill};
  \foreach \x in {0,1,2,3}
   \draw[upfill!45] (\x,0.25)--(\x,3.8);
  \draw[gray,dashed,thick] (1,0.35)--(1,1.15)--(2,1.15)--(2,2.55)--(1,2.55)--(1,3.7);
  \foreach \x/\y in {0/0.8,1/1.7,2/2.1,0/3.15}
   \draw[upfill,very thick] (\x,\y)--++(1,0);
  \draw[upfill,fill=white,thick] (1,0.35) circle (3pt);
  \draw[upfill,fill=white,thick] (1,3.7) circle (3pt);
  \node[right,fill=white,inner sep=2pt] at (1.1,0.35) {down $\to$ up};
  \node[right,fill=white,inner sep=2pt] at (1.1,3.7) {up $\to$ down};
  \node[align=center] at (1.5,-0.35) {New incident marks exchange\\two up colors.};
 \end{scope}
\end{tikzpicture}
\caption{The color-counting comparison for a fixed isolated strand.
Time runs upwards; horizontal segments represent instantaneous exchanges,
and the slot coordinates are schematic. Recoloring the prescribed segment
$X$ up allows its jump marks to be removed and its incident holding intervals
to be filled freely. The only remaining color discontinuities are the two
directed endpoint flips. Each cycle meeting a flip has at most one admissible
coloring, which gives the pin-weight bound used in the path comparison.}
\label{fig:surgery}
\end{figure}


In particular, \eqref{eq:isolateddomination} bounds a joint event, not
the strand law conditioned on isolation.  This is precisely what is
needed when excluding rare isolated strands with large displacement.

\begin{proposition}[Good exposures]\label{prop:good}
Under $\PP_P$ there is an event $\mathcal G$ with
\begin{equation}\label{eq:goodprob}
 \PP_P(\mathcal G^c)=o(\beta^{-5/2})
\end{equation}
having the following properties in $B_{12R_*}(x)$, in either direction
of traversal.
\begin{enumerate}
\item A strand makes at most $c_0L$ inter-site steps in a block and at
most $C\log\beta$ in a bin.  The bounds also apply to segments traced
from an intermediate visit to the region, in either time direction.
\item At each block cut, every ball of radius $c_1L$ centered in the
region contains at most five obstacle slots.  Local regularity at
radius $R_*$ about $x$ holds at all times.  In every block, the set of
fully regular times at that radius has length at least half the block.
\item A strand is called isolated in a block and direction if no other
obstacle slot lies within $c_2L$ of its starting site.  Every such
strand stays at distance greater than two from all other obstacles
and has maximal site displacement at most $CL^{3/4}$ in the block.
\end{enumerate}
The constants are fixed in the order $c_0$, then $c_2$, then $c_1$,
with each sufficiently large for the preceding choices.  The event
and all its properties repeat on subsequent copies of the period.
\end{proposition}

\begin{proof}
First control the number of strand steps.  A chronological connected
list of $j$ inter-site marks in an interval of length $u$, anchored at
a prescribed site, has probability at most
\begin{equation}\label{eq:marklist}
 \frac{(Cu)^j}{j!}.
\end{equation}
Indeed, there are at most $C^j$ successive site and slot-edge choices,
and the ordered time simplex has volume $u^j/j!$; apply the factorial
bound in Proposition~\ref{prop:loopinputs}.  The same count works in
reverse order.  It bounds lists anywhere in the interval, so it also
covers a segment beginning at an arbitrary intermediate visit; no
union over real starting times is necessary.  Taking $j=c_0L$ on
blocks, with $c_0$ large, gives $e^{-cL}$ tails.  Taking
$j=C\log\beta$ on bins gives rapidly small tails by the factorial in
\eqref{eq:marklist}.  A polynomial union over sites, blocks, bins, and
both directions proves the first property with rapidly small error.
For this count and the next one enlarge the region to $B_{20R_*}(x)$.

At a fixed block cut and center, six obstacles in a ball of radius
$c_1L$ supply six distinct slot points there.  There are at most
$CL^{18}$ such six-tuples.  Proposition~\ref{prop:snapshot}, followed
by a union over centers and at most $C\beta$ cuts, bounds this failure
probability by
\begin{equation}\label{eq:sixfailure}
 C\beta R_*^3L^{18}\lambda^6
       \le C\beta^{-3+18\delta}
       =C\beta^{-2.91}=o(\beta^{-5/2}).
\end{equation}
Seam permutations preserve the site counts, including multiplicities,
so this estimate holds on either seam side.

If six blocked sites were present in one of the required fixed-radius
balls during a block, one obstacle at each site could be traced back
to the block cut.  The first property would place their six distinct
starting slots in a ball of radius $c_1L$, once $c_1$ is large enough.
This is impossible.  Thus local regularity holds at all times.
For the remaining condition, Proposition~\ref{prop:geometryinputs}
bounds the expected measure of times when distant-component regularity
fails in any block $I$ by $C|I|e^{-cR_*^\xi}$.  Markov's inequality
and a union over blocks show, with rapidly small error, that it fails
on less than half of every block.  This proves the second property.

The step bound also ensures that an isolated strand cannot meet another
obstacle during its block.  Any potential encounter can be traced back
in both strands to their distinct block-start slots, at total distance
at most a fixed multiple of $c_0L$.  Choose $c_2$ larger than that
multiple.  This argument also covers strands entering the region,
by tracing backward from the encounter.

It remains to bound the isolated strand's displacement.  The site
coordinate of the base walk in Lemma~\ref{lem:isolatedpath} jumps at
rate $S$ to each nearest neighbor.  Exponential martingales in each
coordinate give, for $u\le2L$,
\[
 \PP_{\mathrm b}\left(\max_{0\le v\le u}
                   |X(v)-X(0)|_\infty>CL^{3/4}\right)
       \le C e^{-c\sqrt L}.
\]
For example, use an exponential parameter of order $L^{-1/4}$; the
martingale compensator is at most a constant times $uL^{-1/2}$.
On the step event already established, strands beginning in the
region stay interior.  Lemma~\ref{lem:isolatedpath} therefore bounds
the joint event of isolation and excessive displacement.  A polynomial
union over slots, blocks, and directions is rapidly small.  Seams make
no spatial displacement; for a seam block, enumerate also the slots
on its other seam side.  Combining these exceptional sets proves
\eqref{eq:goodprob} and the third property.
\end{proof}

\subsection{Heat propagation at a fixed good exposure}

From this point, all actual kernels kill on $E$, not at $i$.
The estimates in this subsection are deterministic statements about
the fresh walk at an exposure in $\mathcal G$.  The earlier picture
probabilities have already done their work in constructing that event.

\begin{lemma}[Local heat bound]\label{lem:heatbound}
On $\mathcal G$, a fresh-walk interval kernel of elapsed length $u$,
additionally killed on exit from $B_{R_*}(x)$, has maximum entry at most
\begin{equation}\label{eq:heatbound}
 C(1+u)^{-3/2}.
\end{equation}
The bound holds for either time direction, for every starting slice,
with additional spatial killing, and across copies of the period.
\end{lemma}

\begin{proof}
We use Nash's energy-dissipation argument
\cite[pp.~935--936]{Nash1958}.
For a nonnegative input of mass at most one, let $l=\norm{v}_2^2$.
The all-time local regularity in Proposition~\ref{prop:good} permits
\eqref{eq:nash} at every smooth time, so
$l'\le-cl^{5/3}$.  Instantaneous operations are contractions by
Lemma~\ref{lem:kernelenergy}.  Integrating gives
$\norm{v}_2\le C(1+u)^{-3/4}$ for a unit point input.
Split the kernel at its midpoint and express an entry as the inner
product of a forward mass and an adjoint mass, each propagated for
$u/2$.  Their two norm bounds give \eqref{eq:heatbound}.
\end{proof}

\begin{lemma}[Displacement of the fresh walk]\label{lem:localization}
On $\mathcal G$, from any slot above $B_{6R}(x)$ and any starting
slice, a fresh walk run for elapsed time $u\le2\beta$ in either
direction stays within site distance
\begin{equation}\label{eq:localization}
 C(1+u/L)L^{3/4}
\end{equation}
of its starting site, apart from a rapidly small probability uniform
over the fixed good exposure.  The assertion bounds the probability
of excessive displacement before killing.  In particular, adding
killing on exit from $B_{R_*}(x)$ changes these propagations by a
rapidly small amount in total mass.
\end{lemma}

\begin{proof}
Consider first a piece of one block, of length at most $2L$, starting
at any slice and over a site $z$ in $B_{9R}(x)$.  Stop on exiting
$B_L(z)$, including its exit step in the displacement estimate.
Every obstacle strand visiting this ball or its immediate neighbors
traces back at most $c_0L$ steps to a distinct slot at the deterministic
block cut.  The obstacle count in Proposition~\ref{prop:good} thus
shows that at most five strands can be relevant before stopping.
Their instantaneous neighborhoods have uniformly bounded size.

The rate jumps of each site coordinate have bounded rates and unit
increments.  Their drift vanishes away from the neighboring obstacles
and is uniformly bounded elsewhere.  Lemma~\ref{lem:heatbound}, with
the stopping-ball killing, bounds the expected occupation time near
the relevant strands by
\[
 C\int_0^{2L}(1+t)^{-3/2}\,dt\le C.
\]
The expectation is uniform in the starting slot and slice, including
either side of any prescribed operation.

A forced transport occurs only when the fresh walk is at its prescribed
departure hole.  In any bin there are at most $C\log\beta$ relevant
transports, by the bin step bound for the five strands.  Apply
\eqref{eq:heatbound} to the probability of presence just before each
such operation.  Summing over bins, using their lengths between one
and three, bounds the expected number of forced steps before stopping
by
\[
 C\log\beta\sum_{n\ge0}(1+n)^{-3/2}\le C\log\beta.
\]
The same estimate holds from an arbitrary intermediate slice: a
unit-time interval intersects only a bounded number of the original
bins.  These estimates use deterministic transport times at fixed
exposure, with no conditional Poisson assertion.

Let $A$ be the accumulated occupation time near obstacles plus the
number of forced steps.  From every stopping slice and surviving
state, the conditional expected remaining accumulation is at most
$C\log\beta$.  The Markov property and Markov's inequality show
that each further increment of size $2C\log\beta+1$ has conditional
probability at most $1/2$.  Occupation time is continuous and the
transport-count overshoot is at most one.  Restarting at successive
thresholds therefore gives, for each fixed $a>0$,
\[
 \Pr(A>aL^{3/4})\le C_a
        \exp\{-c_aL^{3/4}/\log\beta\}.
\]
This controls both the integrated drift and the forced displacement.

Subtract the integrated drift from the ordinary rate jumps.  For each
coordinate, bounded rates and unit jumps imply that its compensated
process $M_t$ has
$\exp\{\theta M_t-C\theta^2t\}$ as a supermartingale for
$|\theta|\le1$, with a fixed $C$.  The maximal inequality with
$\theta$ of order $L^{-1/4}$ gives
\[
 \Pr\bigl(\max_{t\le2L}|M_t|>aL^{3/4}\bigr)
                                  \le C_a e^{-c_a\sqrt L}.
\]
The three displacement contributions prove the claimed bound on one
block piece.  Since $CL^{3/4}<L$ for large $\beta$, they also bound
the probability of reaching the stopping sphere.  Seam operations
change no site coordinate.

An interval of length $u$ meets at most $C(1+u/L)$ block pieces.
Apply the preceding uniform bound successively and add their
displacements.  There are polynomially many pieces, so the total
failure probability remains rapidly small.  Throughout an interval
of length at most $2\beta$, the resulting displacement is
$O(\beta L^{-1/4})=o(R)$.  The walk therefore stays in the local
region used above, and in $B_{R_*}(x)$.  This proves the lemma.
\end{proof}

\subsection{Overlaps of flows}

The final geometric estimate measures the energy cost of superposing
flows from different source locations.  It retains their separation,
which will be essential when averaging the sparse source locations.

\begin{lemma}[Flow overlap]\label{lem:flowoverlap}
At a fully regular slice on $\mathcal G$, choose the unit flows in
\eqref{eq:flowenvelopes} for available slots $a,a'$ over $B_{6R}(x)$.
Then
\begin{equation}\label{eq:flowoverlap}
 \sum_e|J_a(e)J_{a'}(e)|
       \le\frac{C}{1+\dist(a,a')},
\end{equation}
where the distance is between their sites.
\end{lemma}

\begin{proof}
Where both edge distances from the source sites are at most $R_*$,
use the two inverse-square envelopes.  Their convolution in dimension
three is bounded by $C/(1+b)$ at source separation $b$: balls of
radius $b/2$ about either source contribute $C/(1+b)$, and dyadic
shells outside them have the same total bound.

If one source distance exceeds $R_*$, both are at least a fixed
multiple of $R_*$, since $a,a'$ lie over $B_{6R}(x)$ and $R=o(R_*)$.
The global envelopes then bound the remaining sum by
\[
 C R_*^{3-2(199/100)}
       =C R_*^{-98/100}
       \le C R^{-1}.
\]
The last inequality follows from
$(5/3)(98/100)=49/30>3/2+\delta/4$.
As $b\le12R$, this is bounded by the right-hand side of
\eqref{eq:flowoverlap}.
\end{proof}
